\cotefile{157-1}
\scanpage{227}\mspage{82}
\[
\Phi\longleftrightarrow\Psi\quad\Longleftrightarrow\quad
\varphi^{-1}(\Phi\times\Psi)\subset\unsure{\mathcal{D}|\mathscr{D}}_0
\]
Mais à vrai dire, je m'aperçois que c'est
quand $=$ \unsure{du dual} d'une relation \deleted{\unsure{entre} $\operatorname{Fl}($}
\unsure{il|qu'il} \deleted{\unsure{s'agit}} \unsure{en|sur} $\operatorname{Fl}(M)$, \deleted{\unsure{puisque}} et cette relation
\deleted{\unsure{peut s'é}} entre $i:A\unsure{-|\to}B$ et $p:X\to Y$
peut s'écrire \unsure{ainsi|comme} :
\[
\begin{aligned}
&\text{pour ttes flèches } A-X,\ B-Y\text{ rendant}\\
&\text{commutatif le diagramme carré, la}\\
&\text{flèche }\varphi\text{ pointillée du carré existe.}
\end{aligned}\tag*{$(**)$}
\]

Cela me fait penser que si on a
une relation $\mathcal{R}$ entre $\unsure{\mathfrak{P}|\mathcal{P}}(P)$ et $\mathfrak{P}(Q)$,
satisfaisant $\unsure{D_2|D|D_1}$, donc décrite par une
application décroissante
\[
\mathfrak{P}(P)\longrightarrow\mathfrak{P}(Q)
\]
\unsure{transformant} $\operatorname{Sup}$ en $\operatorname{Inf}$, \deleted{\unsure{telle}} \inserted{la donnée d'}une
telle application équivaut à celle d'une
application quelconque
\[
P\unsure{-|\to}\mathfrak{P}(Q),
\]
ou ce qui revient au $=$, d'une relation
$R$ sur $P\times Q$. Donc \unsure{cette} structure
\deleted{\unsure{devait}} \unsure{doit|peut} être regardée comme ``de première
espèce''. La situation change à partir
du moment où, pour décrire la relation
$\mathcal{R}$,\inserted{ ou les propriétés voulues,} on se trouve \unsure{contraint} de remplacer
$\mathfrak{P}(P)$, $\mathfrak{P}(Q)$ par des sous-ensembles
convenables -- quand on ne regarde que
\uline{\unsure{certaines}} parties de $P,Q$. C'est ce qu'on
% archive sheet number: 227
\scanpage{228}\mspage{83}
va faire au § suivant. Le cas que je vais y
traiter est le seul de son espèce que j'aie
\deleted{\unsure{co}}rencontré -- un cas de ``dualité formelle'' qui
s'introduit de façon naturelle, et qui
\unsure{ne} soit ``de $1^{\text{\unsure{ière|ère}}}$ espèce''.

6. \uline{Transporteurs}, \uline{noyaux}, \uline{stabiliseurs relatifs à} $W$.

Soit $M$ une catégorie, munie d'une partie
\[
W\subset\operatorname{Fl}(M).\tag*{(6.1)}
\]
Je pense au cas d'un localiseur, mais n'introduirai d'hypothèse sur $W$ qu'au besoin.
Pour simplifier, je suppose $M$ stable
par produits fibrés.

Je pose \margin{cf rectification p.~86}
\[
\left\{\begin{aligned}
P&=\unsure{\mathfrak{P}|\mathcal{P}}(W)=\{A\in\mathfrak{P}(\operatorname{Fl}(M))\mid A\subset W\}\\
Q&=\{L\in\mathfrak{P}(\operatorname{Fl}(M))\mid L\ \text{\uline{stable} par chgt de}\ \}\\[-2pt]
 &\qquad\text{base }\unsure{(}\text{par quoi je sous-entends}\\[-2pt]
 &\qquad\text{aussi, stable par}\\[-2pt]
 &\qquad\text{isomorphie) et par}\\[-2pt]
 &\qquad\text{composition}.
\end{aligned}\right.\tag*{(6.2)}
\]

J'introduis une relation
\[
\mathcal{R}\subset P\times Q
\]
de \unsure{fa|la|\deleted{fa}} façon suivante. Pour $A\in P$, $L\in Q$,
\unsure{donc|on} pose
\diagram{\begin{tikzcd}[ampersand replacement=\&,column sep=small,row sep=small]
{} \& {} \& {\text{Pour \unsure{tout|tt} diagramme cartésien}} \& {} \\
{\text{(6.3)}\quad(A,L)\in\mathcal R\ \Longleftrightarrow} \& X\arrow[d,"f"] \& X'\arrow[l]\arrow[d,"\unsure{f'|g'}"] \& {\text{où }f\in A,\ p\in L,}\\
{} \& Y \& Y'\arrow[l,"p"] \& {\text{on a } f'\in\deleted{\unsure{A}}W}
\end{tikzcd}}
(i.e. \deleted{\unsure{le chgt de base par un} $p\in L$}
\inserted{$A$ est stable par chgt de base}
\deleted{\unsure{d'un}} \inserted{par} $\unsure{p|f}\in\unsure{L|A}$)
\deleted{\unsure{donne un} $f'\in W$)}.
% archive sheet number: 228
\scanpage{229}\mspage{84}
Quand j'ordonne \deleted{$\unsure{P=\mathfrak{P}(W)}$} $\unsure{P}, Q\subset\unsure{\mathfrak{P}|\mathcal{P}}(\operatorname{Fl}(M))$
par la relation d'ordre induite par l'inclusion
de parties de $\operatorname{Fl}(M)$, on trouve des \inserted{\textsuperscript{ss}/}ensembles
ordonnés de $\mathfrak{P}(\operatorname{Fl}(M))$, stables par réunions
quelconques, et aussi par intersections quelconques
avec la seule réserve que dans le cas de
$P$, la famille \deleted{\unsure{doit}} d'indices doit être $\unsure{\ne|=}\varnothing$
(puisqu'alors on trouve comme intersection
$\operatorname{Fl}(M)$ \deleted{tout} entier, qui n'est pas en général
\unsure{contenue|contenu} dans $W$).

On va vérifier la condition $(\unsure{D_0|D})$ de la
page 69, pour $\mathcal{R}$ et les relations d'ordre.
Il y a deux vérifications à faire

$1^\circ)$ Donnons nous $A\in P$\inserted{, i.e. $A\subset W$}, et considérons l'ensemble
$\mathcal R(A)$ des $L\in Q$ tels que l'on ait \deleted{$\unsure{R\in}$}
$(A,L)\in\mathcal R$. Il est clair que c'est un ouvert
(et de $=$, que $\,{}^{\unsure{c|t}}\mathcal R(L)$ est un ouvert \inserted{de $\unsure{P|Q}$} pour
tout $L$ dans $Q$) -- i.e. $L\in R(A)$, $L'\subset L$
implique $L'\in R(A)$. Reste à trouver un
plus grand élément de $\mathcal R(A)$. Notons que
si $L\in R(A)$, alors \deleted{le chgt de base}\inserted{pour tout}
$p\in L$, $p:Y'\unsure{-|\to}Y$, le changt de base
$p^{*}:\operatorname{Cat}_{/Y}\unsure{-|\to}\operatorname{Cat}_{/Y'}\ :\quad X\unsure{\vert\!-\!\vert|\mapsto|\leftarrow}X\times_Y Y'$
doit transformer flèches $f\in A$ en flèches $g\in W$,
\diagram{\begin{tikzcd}[ampersand replacement=\&,column sep=small,row sep=small]
X_1\arrow[d,"f"'] \& X'_1\arrow[l,"\unsure{q_1|p_1}"']\arrow[d,"g"]\\
X_2\arrow[d] \& X'_2\arrow[l,"\unsure{q_2|p_2}"']\arrow[d]\\
Y \& Y'\arrow[l,"p"']
\end{tikzcd}}
comme il résulte du fait que
$L$ est stable par chgt de base.
Donc $L\subset\unsure{\mathbb{L}|L}_W(A)$, où on pose
% archive sheet number: 229
\scanpage{230}\mspage{85}
\[
\begin{aligned}
\overset{\smash{\tikz{\draw (0,0)--(0,.08)--(.95,.08)--(.95,0);}}}{\mathbb L_{\deleted{\unsure{w}}W}(A)}
 &=\{(p:Y'\to Y)\in\deleted{\unsure{\operatorname{Fl}}}\operatorname{Fl}(M)
 \mid p^*:\operatorname{Cat}_{/Y}\unsure{-|\to}\operatorname{Cat}_{/Y'}\ \text{transforme}\\
 &\hspace{6em}\text{flèches de }A\text{ en flèches \uuline{de $W$}}\}.
\end{aligned}\tag*{(6.4)}
\]
\diagram{\begin{tikzcd}[ampersand replacement=\&,column sep=1.2em,row sep=1em]
X_1\arrow[d,"f"'] \& X'_1\arrow[l]\arrow[d,"f'"]\\
X_2\arrow[d]\arrow[ur,phantom,"\mathrm{cart}"] \& X'_2\arrow[l]\arrow[d]\\
Y\arrow[ur,phantom,"\mathrm{cart}"] \& Y'\arrow[l,"p"]
\end{tikzcd}}
\margin{$f\in A\Rightarrow f'\in W$ \deleted{\unsure{$f'$}}}
\margin{NB la \unsure{clause peut être|clause précédente} \unsure{lue} en termes de diagr. cartésiens, ce qui \unsure{en est} en fait \unsure{indép. de} l'hypothèse de stabilité de $W$ par isoms.}

Il est clair d'autre part que cette partie de $\operatorname{Fl}(M)$
est stable par changt de base\inserted{, \deleted{et par composition}}, \unsure{donc} $\mathbb L_W(A)\in Q$,
\diagram{\begin{tikzcd}[ampersand replacement=\&,row sep=1em,column sep=small]
X_1\arrow[d,"f"'] \& X'_1\arrow[l,"q_1"']\arrow[d,"f'"]\\
X_2\arrow[d]\arrow[ur,phantom,"\mathrm{cart}"] \& X'_2\arrow[l,"q_2"']\arrow[d]\\
Y_1\arrow[d]\arrow[ur,phantom,"\mathrm{cart}"] \& Y'_1\arrow[l,"p_1"']\arrow[d]\\
Y\arrow[ur,phantom,"\mathrm{cart}"] \& Y'\arrow[l,"p"]
\end{tikzcd}}
et que $(A,\mathbb L_W(A))\in R$. Donc
$\mathbb L_W(A)$ est un plus grand élément
de $\mathcal R(A)$, donc on trouve
\[
\mathcal R=\{(A,L)\in P\unsure{\times|<}Q\mid L\subset\mathbb L_W(A)\}.\tag*{(6.5)}
\]

Inversement, partons de $L\in Q$,
et montrons que l'ouvert $\,{}^{t}\mathcal R(L)$ de $P$ a
un plus grand élément. Pour ceci,
posons
\diagram{\begin{tikzcd}[ampersand replacement=\&,column sep=0.7em,row sep=0.7em]
{} \& {} \& X\arrow[d,"f"'] \& X'\arrow[l,"q"']\arrow[d,"f'"]\\
{\text{(6.6)}\quad\unsure{\mathbb A|A}_W(L)=\big\{(f\unsure{\in|:}X\to Y)\in\inserted{W\subset}\operatorname{Fl}(M)\ \big|} \& {\forall\,\deleted{\unsure{p}}\,:} \& Y\arrow[ur,phantom,"\text{\uuline{cart}}"] \& {Y',}\arrow[l,leftrightarrow,"{\unsure{p}}"]\\
{} \& {} \& {} \& {\text{avec}}\\
{} \& {} \& {} \& {p\in L,\ \deleted{\text{\unsure{le changt}}}}\\
{} \& {} \& {} \& {\deleted{\text{\unsure{induit}}}\inserted{\text{on a}}\,f'\in W\big\}.}
\end{tikzcd}}

(NB Quand $L$ contient les identités, comme c'est
le cas notamment s'il est de la forme $\mathbb L_W(A)$,
-- alors il contient les iso, \unsure{du moins} si $A$ les
contient -- alors $\mathbb A_W(L)$ est l'ens des $f\in\operatorname{Fl}(M)$
satisfaisant à la condition énoncée après $\mid$, sans
avoir à spécifier $f\in W$.) On a $\mathbb A_W(L)\in P$
\deleted{\unsure{par} $A$} et \deleted{il est clair que}\inserted{je vérifie que} $(\mathbb A_W(L),L)\in\mathcal R$,
\inserted{(en utilisant la stabilité de $L$ par composition)}.
\unsure{enfin} $(A,L)\in\mathcal R$ implique $A\subset\mathbb A_W(L)$, \unsure{(ce}
\inserted{qui est clair, en utilisant que $A\subset W$)},
donc $\mathbb A_W(L)$ est un plus grand élément de
$\,{}^{t}\mathcal R(L)$.
% archive sheet number: 230
\scanpage{231}\mspage{86}
\deleted{À dire vrai} Notons que si \uline{$W$ est stable par composition} (ce que
nous allons supposer désormais), il en est de même
de $\mathbb A_W(L)$ (ce qui utilise \deleted{de façon} encore
le fait que $L$ soit stable par changt de base).
De plus, si \uline{$W$ contient les iso} (ce que nous
supposerons également), $\mathbb A_W(L)$ les contient aussi.
Quant à $\mathbb L_W(A)$, il contient les iso.
\uline{pourvu que} $A$ soit stable par isomorphisme\unsure{;}
quand $A$ est stable par composition et contient les
\unsure{iso|id}, cette condition sera automatiquement satisfaite. Nous tenons à cette condition, c'est
pourquoi du côté de $A$, dans la définition
de $P$, \unsure{je} vais introduire cette condition
supplémentaire que $A$ soit \uline{stable par iso}

Donc je pose, en corrigeant (6.2)
\[
\left\{\begin{aligned}
P={}&\{A\in\unsure{\mathfrak{P}|\mathcal{P}}(\operatorname{Fl}(M))\mid
\deleted{A\subset W,\ \unsure{A}\ A\ \text{stable par iso}\}}\\
&\inserted{\operatorname{Is}(M)\subset A\subset W,\ A\ \text{stable par composition}}\},\\
\deleted{\text{\unsure{Notons que}}}\quad
Q={}&\{L\in\mathfrak{P}(\operatorname{Fl}(M))\mid\operatorname{Is}(M)\subset L,\ L\ \text{stable }\deleted{\text{par}}\\
&\qquad\deleted{\text{composition et}}\ \text{par chgt. de base}\}.
\end{aligned}\right.\tag*{(6.7)}
\]

On est bien sûr intéressé surtout aux
\unsure{éléments|élémts} clos de $P,Q$ i.e. les parties
$A,L$ de $\operatorname{Fl}(M)$, \unsure{contnt} les iso, stables
par composition, (et $L$ stable aussi par chgt
de base, $A\subset W$) telles que l'on ait
\[
A=\mathbb A_W(L),\quad L=\mathbb L(A)\tag*{(6.8)}
\]
% archive sheet number: 231
\scanpage{232}\mspage{87}
Je garderai les notations $\unsure{\sharp|\natural},\flat$ \unsure{du|des} \unsure{par.} \unsure{précédent|précédents}.
Quels sont les deux couples extrêmes ?

Partant de $A_0=\operatorname{Is}(M)$ (plus petit des éléments
de $P$), on trouve $\mathbb L_W(A_0)=\operatorname{Fl}(M)$ (plus grand
élément de $Q$), et $A_0^\sharp=\mathbb A_W(\operatorname{Fl}(M))=W_{\mathrm{univ.}}$,
ensemble des flèches qui sont dans $W$ et le
restent par tt chgt de base :
\[
\deleted{\operatorname{Is}(M)^\sharp=W_{\mathrm{univ}},\quad\operatorname{Is}(M)}
\]
\[
A_0^\sharp=W_{\mathrm{univ}},\quad
\deleted{\unsure{A|\alpha}}\beta_0(A_0)\overset{\text{déf}}{=}\mathbb L(A_0)=\operatorname{Fl}(M),
\tag*{(6.9.)}
\]
\margin{où $A_0=\operatorname{Is}(M)$, plus petit él. de $P$}

Inversement, partons de $L_0=\operatorname{Is}(M)$, le plus petit
élément de $Q$. On a $\mathbb A_W(L_0)=W$ ($W$
étant stable par composition et \unsure{contnt} les iso),
et $L_0^\flat=\mathbb L(W)=\operatorname{Fib}_W$, les $W$-fibrations.
\[
L_0^\flat=\operatorname{Fib}_W,\quad \alpha_0(L_0)=W,
\tag*{(6.10)}
\]
\margin{où $L_0=\operatorname{Is}(M)$, plus petit él. de $Q$.}

Donc les parties $\mathcal R$-closes à g. de $\operatorname{Fl} M$
se promènent entre $W_{\mathrm{univ}}$ et $W$, les
parties $\mathcal R$-closes à droite entre $\operatorname{Fib}_W$ et $\operatorname{Fl} M$.

\uline{Proposition} 6.1. Soient $A\in P$, $L\in Q$ tels
\llap{\smash{\rule[-4.9\baselineskip]{0.35pt}{5.2\baselineskip}}\hspace{0.5em}}que $A=\deleted{\unsure{\beta}}\alpha_0(L)$, $L=\beta_0(A)$ (i.e. $A=\mathbb A_W(L)$,
$L=\mathbb L_W(A)$). Pour que $L$ soit stable par
composition, il f. et s. que $A$ soit
stable par chgt. de base \deleted{pour} \unsure{par les|les}
flèches $f\in L$. (Plus généralement si $L\inserted{\in Q}$ est
un élément quelconque de $Q$, et si on suppose de
% archive sheet number: 232
\scanpage{233}\mspage{88}
de plus que \deleted{$\unsure{A}$} $L$ est stable par composition,
alors $A=\mathbb A_W(L)$ est stable par chgt de base
par $L$.)

Terminologie : pour $A\in P$ i.e. $A$ une
partie de $\operatorname{Fl} M$ \deleted{\unsure{comprenant}}\inserted{contenue} dans $W$,
contenant les $\operatorname{Is}$ et stable par composition,
$\mathbb L_W(A)$ s'appelle le \uline{transporteur} \uline{de} $A$ \uline{dans} $W$.
(Il y a une notion duale de \uline{cotransporteur}.)

C'est \inserted{un \unsure{élém|élément} de $Q$, donc} une partie de $\operatorname{Fl} M$ contenant les
iso et stable par \deleted{composition}\inserted{chgt. de base}. D'autre
part, pour $L$ dans $Q$ \deleted{\inserted{\unsure{i.e. $L$ une}} (partie de $\operatorname{Fl}(M)$,
contenant les iso et stable par chgt de base)},
les $f\in\mathbb A_W(L)$ s'appellent flèches \uline{$W$-\deleted{\unsure{L}}-\unsure{asphériques}},
ou \uline{$W$-$L$-équivalences}, \unsure{ou les|ou} \uline{$W$-équiv. $L$-universelles}\unsure{?|).}
\unsure{Et} $\mathbb A_W(L)$ s'appelle le \uline{$L$-noyau de} $W$.

On pourrait le noter plus simplement $W_L$
(quand cette notation ne \unsure{conflicte} pas avec
$W_S\subset\operatorname{Fl}(M_{/S})$ pour $S\in\operatorname{Ob} M$). C'est une
partie de $W$ qui (comme $W$ lui-$=$)
contient les iso et est stable par composition.

Enfin, il y a lieu aussi d'introduire
le \uline{stabiliseur} $\operatorname{Stab}(A)$, \deleted{\unsure{qui}} transporteur de
$A$ dans $A$, qui ne dépend pas de la donnée
de $W$, mais seulement de la partie $A$ de
$\operatorname{Fl} M$, contenant les iso et stable par composition.
\deleted{Il \unsure{pourra}} On pourrait aussi le noter $\operatorname{Fib}_A$, et
appeler ses éléments des \uline{$A$-fibrations} ($A$ jouant à
présent le rôle de $W$), \deleted{Il semblerait que,}
On a $\operatorname{Stab}(A)\subset\mathbb L_W(A)$, et il semblerait qu'en
% archive sheet number: 233
\scanpage{234}\mspage{89}
pratique, ce sont \unsure{les|ces} stabiliseurs, plus que les
transporteurs $\mathbb L_W(A)$, qui sont importants (quand
les deux ne coïncident pas). C'est ainsi que
\inserted{dans $\operatorname{Cat}$, si $W$ est un local. fond.,}
les morphismes $W$-lisses (resp\textsuperscript{\unsure{t}} $W$-propres)
$\operatorname{Liss}_W$, $(\operatorname{Prop}_W)$, qui sont les stabiliseurs des
ensembles $W_{\mathrm{as}}$ $(W_{\mathrm{cas}})$ des \unsure{ens} de morphismes
$W$-asphériques (resp. $W$-coasphériques),
jouent un rôle essentiel alors que je ne
sais pas caractériser de façon sensible les
flèches des transporteurs $\mathbb L_W(W_{\mathrm{as}})$, $\mathbb L_W(W_{\mathrm{cas}})$,
lesquels contiennent, en plus des morphismes
$W$-lisses (resp. $W$-propres) \deleted{les} $W$-fibrations
$\operatorname{Fib}_W$. Je ne sais pas \inserted{comment} caractériser une
classe de morphismes qui apparaît aussi
hétéroclite ! De même, il y a lieu
d'introduire et de caractériser le
stabiliseur de $W_{\mathrm{bias}}$ (ens. des flèches
$W$-biasphériques), comme les morphismes
``$W$-nets'' $\operatorname{Net}_W$, \deleted{\unsure{par}} une description
similaire à celle des $\operatorname{Liss}_W$, $\operatorname{Prop}_W$,
utilisant les bilocalisés $\underset{u}{{}_{\backslash}X_{/}}$, $\underset{\unsure{v|u}}{{}_{\backslash}Y_{/}}$ au lieu
des localisés $X_{/x},Y_{/y}$ ou des colocalisés
$\,{}_{x\backslash}X,{}_{y\backslash}Y$. Il y a lieu enfin de déterminer
le stabiliseur de l'ensemble $\unsure{\operatorname{Equ}|\operatorname{Eqv}}$ des
$\operatorname{Cat}$-\unsure{équivalences} (ce qui ne dépend pas de la
donnée d'un $W$\unsure{)}, -- mais on a $\operatorname{Equ}\subset\unsure{\mathfrak{P}|\mathcal{P}}(W)$
\inserted{pour tt $W$ fondamental}).
% archive sheet number: 234
\scanpage{235}\mspage{90}
7. \uline{Morphismes cofiltrants et loc cofiltrants}.

Mais je ne veux pas m'attarder à présent sur
ces développements, qui ne me paraissent pas
urgents, et je voudrais revenir aux questions
liées à la notion de catégories filtrantes \deleted{\unsure{et}}
cofiltrantes.

% slanted
\margin{J'ai \unsure{modifié} la terminologie, \unsure{maintenant} \unsure{cof}, loc. fbt cof, \unsure{pré} \illegible{1} ; fbt cof = \unsure{$0$-connexité|$0$-conn.} \unsure{stable par passage|stable par prod. fibrés} : \unsure{si|tt} $X'\unsure{-|\to}\unsure{X|x}$ \unsure{dans} $X$ cof, $X'$ étale et \unsure{connexe}.}
% slanted
\margin{Revoir la terminologie du cas (loc. fbt cofiltrant, \unsure{etc.}), qui est une situation très \unsure{ingrate}, faute d'être stable par passage \unsure{aux|un} $X^{\unsure{\prime|\prime\prime}}$ \unsure{étal|étale} sur $X$.}
\uline{Définition} 7.1. Soit $f:X\to Y$ une flèche dans $\operatorname{Cat}$.
On dit que $f$ \uline{est cofiltrant} (resp. \uline{loc. cofiltrant},
\uline{faiblement cofiltrant}$^{*}$, respt \uline{loc. faiblement cofiltrant}
-- ou \uline{précofiltrant}) si pour tt $y$ dans
$Y$, la catégorie \inserted{colocalisée} $\,{}_{y\backslash}X$ est cofiltrante (resp. \ldots)

[\inserted{Dualmt,} On dit que $f$ est filtrant, resp. \ldots, si les
catégories $X_{/y}$ sont filtrantes, resp. \ldots, i.e.
si $f^\circ:X^\circ\unsure{-|\to}Y^\circ$ est cofiltrant resp \ldots]

\deleted{\unsure{Proposition|Remarque}} \uline{\unsure{Remarque} 7.2.} \unsure{Notions} stables par chgt de base coétale
\inserted{contient les iso, \unsure{et|de} $=$ \unsure{bien sûr|les quatre genres} \unsure{imm}}\unsure{ersions}
% slanted
\margin{\deleted{\illegible{1}} et \unsure{les} \unsure{immersions} \unsure{coétales} sont \unsure{plates}.}
\uline{Proposition} 7.3. Soit $f:X\unsure{-|\to}Y$ une \deleted{flèche}
cofiltrant (resp. loc. cofiltrant, \unsure{respt.} fbt cofiltrant,
resp. \deleted{loc. fbt \unsure{pr}} précofiltrant), et supposons
que $Y$ soit cofiltrant (resp. loc cofiltrant,
resp fbt cofiltrant, resp. précofiltrant). Alors
$X$ lui-même est cofiltrant (resp. \ldots).

\uline{Dém}. Il s'agit, dans les quatre cas
envisagés, de vérifier pour $X$ la condition
\[
\underbrace{\mathrm{PS}_g1'\ /}_{\text{\unsure{deux autres cas|deux premiers cas}}}
\ \mathrm{PS}_g1\ /\ \mathrm{PS}_g1'+\mathrm{PS}_g2\ /\ \mathrm{PS}_g1+\mathrm{PS}_{\unsure{g|s}}2.
\]
Mais
% archive sheet number: 235
\scanpage{236}\mspage{91}
pour chacune des trois conditions $\mathrm{PS}_g1$, $\mathrm{PS}_g1'$, $\mathrm{PS}2$
séparément, on voit immédiatement que si
elle est vérifiée pour $Y$ et pour les $\,{}_{y\backslash}X$,
elle l'est pour $X$, d'où la proposition.

\uline{Corollaire} 7.3. Les quatre notions envisagées
dans 7.1. sont \deleted{\unsure{filtrantes}}\inserted{stables} par composition.

\diagram{\begin{tikzcd}[ampersand replacement=\&,row sep=small,column sep=small]
X\arrow[d] \& {}_{z\backslash}X\arrow[l]\arrow[d]\\
Y\arrow[d] \& {}_{z\backslash}Y\arrow[l]\arrow[d]\\
Z \& {}_{z\backslash}Z\arrow[l]
\end{tikzcd}}
Démonstration immédiate, en
utilisant la remarque 7.2.

\uline{Remarque} \unsure{7.4.|?.4.} 1) Notons que
la notion de morphisme cofiltrant
\inserted{ou fbt cofiltrant}
\deleted{(et les variantes 7.1)} n'est \uline{pas} stable par
chgt. de base quelconque, ni $=$ par une
\unsure{imm.} ouverte. Exemple : $e\unsure{-|\to}\Delta^1$, appliquant
$e$ en $1$, est cofiltrant, mais sa
fibre en $0$ est vide donc n'est
ni cofiltrante, ni fbt. cofiltrante.
\diagram{\begin{tikzcd}[ampersand replacement=\&,row sep=0.7em,column sep=small]
{} \& {\llap{$\bullet$\hspace{0.8em}}}\arrow[d,mapsto]\\
0\arrow[r] \& 1
\end{tikzcd}}

Je n'ai pas \unsure{vu|eu} d'exemple montrant que la
notion de flèches \uline{loc}. cofiltrantes resp. précofiltrantes
n'est pas stable par changt. de base, et ne voudrais pas m'y attarder, n'ayant pas besoin de
connaître la réponse.

2) Pour que $X$ soit cofiltrante (resp. \ldots),
il f. et s. que $X\longrightarrow e$ le soit. Ainsi la
notion absolue est ramenée à une notion
relative.

3) Si $f:X\to Y$ est cofiltrante, elle est
% archive sheet number: 236
\scanpage{237}\mspage{92}
$W$-coasphérique (donc $\in W$) pour tt \deleted{\unsure{fonct}} $W$
satisfaisant l'axiome $\mathrm{Loc}\ 8$ des limites
inductives. En particulier, elle est $W_0$-\unsure{coasph},
i.e. un foncteur cofiltrant \deleted{\unsure{ou}} (et $=$
seulmt fbt cofiltrant) est \deleted{\illegible{1}} \uline{cofinal}.

La raison principale pour m'intéresser
aux flèches cofiltrantes est contenue
dans la

\uline{Proposition} 7.4. Soit $f:X\longrightarrow Y$ \inserted{dans $\operatorname{Cat}$,} loc. cofiltrant
(resp. loc. \uline{fbt.} cofiltrant)
Alors : \deleted{pour toute}

a) Pour tout $X'$ étale sur $X$, le composé $f':$
$X'\unsure{-|\to}X-Y$ est loc. cofiltrant\inserted{ (resp. \ldots)}, et pour
toute factorisation de $f'$ en $X'\xrightarrow{g'}Y'\longrightarrow Y$,
avec $Y'-Y$ étale, $g'$ est loc cofiltrant (resp. \ldots).

b) \deleted{\unsure{Les parties libres}}, \inserted{$\forall x$ dans $X$,} le \unsure{morphisme} sur les localisés
\deleted{\unsure{sont}} $f_{/x}:X_{/x}\unsure{-|\to}Y_{/y}\qquad(y=f(x))$
est \uline{cofiltrant} (resp. fbt cofiltrant),
\deleted{donc $\unsure{W_0}$-coasphérique, i.e. \uline{cofinal}.}

\deleted{Pour} c) Si $x$, \deleted{\unsure{est un}} comme dans b), et
$U$ est un ouvert de $X_{/x}$, désignant
par $V$ l'ouvert de $Y_{/y}$ engendré par
l'image de $f_{/x}$ (i.e. l'ouvert formé des
$z$ dans $Y_{/y}$ tels que $\,{}_{z\backslash}(X_{/x})$ soit non vide),
la flèche induite $g:U\longrightarrow V$ \unsure{est|\deleted{est}} \inserted{cofiltrante (resp. fbt \unsure{cofiltrte}), \unsure{a fortiori}} \uline{cofinale},
(et $\unsure{\overline{m}|=}$ $W$-coasphérique, pour toute $W$
satisfaisant l'axiome des limites $\mathrm{Loc}\ 8$).
% slanted
\margin{Il suffit ici qu'on soit dans l'hyp. \unsure{resp}\textsuperscript{\unsure{ve}} \deleted{plus} ($f$ \uline{fbt}. loc. cofilt\unsure{(ant)}) \unsure{Sauf} pour le \unsure{NB ici}.}
% archive sheet number: 237
\scanpage{238}\mspage{\unsure{93}}
Donc pour tt \uline{co}faisceau \inserted{$\unsure{E|\mathcal E}$} sur $V$, à valeurs dans
une catégorie $M$ avec $\underset{\unsure{-|\longrightarrow}}{\lim}$, on a iso
\[
\underset{\unsure{U|u}}{\varinjlim}\,g^*(E)\stackrel{\sim}{\longrightarrow}
\underset{V}{\underset{\unsure{-|\longrightarrow}}{\lim}}\,E,
\]
et pour tt préfaisceau \inserted{$F$} sur $V$, à val. dans une
catégorie $M$ avec $\varprojlim$, on a iso
\[
\underset{U}{\varprojlim}\,g^*(F)\stackrel{\sim}{\longleftarrow}
\underset{V}{\varprojlim}\,F.
\]

Démonstration. \deleted{\unsure{Comme}} \inserted{Je vais \deleted{\unsure{ici}} d'abord expliciter}
\deleted{les morphismes étales sont loc. cofiltrants}

\uline{Lemme} \inserted{7.5.} Soit $h:X'\unsure{-|\to}X$ un morphisme étale
dans $\operatorname{Cat}$. Alors $h$ est loc. cofiltrant. \deleted{\unsure{car en}}\inserted{De façon}
plus précise, pour tout $x$ dans $X$, \deleted{$\,{}_{x\backslash}X'$ est
\unsure{composé} \inserted{\unsure{somme}} \unsure{directe} de ses composantes connexes,}
on a
\[
\pi_0({}_{x\backslash}X')\overset{\sim}{\longrightarrow}\operatorname{Ob} X'_x
\]
toute composante connexe de $\,{}_{x\backslash}X'$ contient
un élément et un seul de \deleted{$\unsure{X'|X'_x}$} la
cat. discrète $X'_x\subset{}_{x\backslash}X'$. Si $x'$ est dans
$X'_x$ (i.e. \deleted{\unsure{avec}} $x'$ objet de $X'$ tel que $h(x')=x$),
c'est un objet \uline{initial} de la composante
connexe qui le contient. (Donc cette composante est cofiltrante.)

Ceci \unsure{prouve}\deleted{\unsure{donc}} \deleted{\unsure{la (ad}} (\unsure{admis cela|admet cela}) la première
assertion de 7.4., compte tenu de 7.3.
(stabilité par composition), on \unsure{prouve|perçoit} qu'une
% archive sheet number: 238
\scanpage{239}\mspage{94}
\diagram{\begin{tikzcd}[ampersand replacement=\&,column sep=2em,row sep=1.4em]
{}_{y\backslash}X'\arrow[r]\arrow[d,"\text{ét}"']\arrow[dr,phantom,"\text{cart}"] \& X'\arrow[d,"\text{ét.}"]\\
{}_{y\backslash}X\arrow[r]\arrow[d]\arrow[dr,phantom,"\text{cart}"] \& X\arrow[d]\\
{}_{y\backslash}Y\arrow[r] \& Y
\end{tikzcd}}
catégorie étale (\unsure{en} l'occurrence $\,{}_{y\backslash}X'$)
au dessus d'une catégorie
loc. cofiltrante (resp loc fbt. cofiltrante) (en l'occurrence $\,{}_{y\backslash}X$) est
itou. Pour la deuxième assertion de a), on
peut \unsure{maint}\deleted{\unsure{enant}} supposer $X'=X$. Il faut prouver
que pour tout $y'$ de $Y'$,
\diagram{\begin{tikzcd}[ampersand replacement=\&,column sep=2.3em,row sep=1.6em]
{}_{y'\backslash}\unsure{Y|X}\arrow[r]\arrow[d,"\text{ét}"']\arrow[dr,phantom,"\text{\uline{cart}}"] \& {}_{y\backslash}X\arrow[r,dash]\arrow[d,dash,"\text{ét}"]\arrow[dr,phantom,"\text{\uline{cart}}"] \& X^{\unsure{g|\prime}}\arrow[d,dash,"\text{ét}"]\\
{}_{y'\backslash}Y'\arrow[r] \& {}_{y\backslash}Y'\arrow[r]\arrow[d]\arrow[dr,phantom,"\text{\uline{cart}}"] \& Y'\arrow[d]\\
\& {}_{y\backslash}Y\arrow[r] \& Y
\end{tikzcd}}
$\,{}_{y'\backslash}X\mathrel{\overset{\mathrm{déf}}{=}}X\times_{Y'}{}_{y'\backslash}Y'$ est loc. cofiltrante
(resp loc. fbt. cofiltrante). \deleted{\unsure{Soit}}
$y$\inserted{ \unsure{en $Y$|en $y$}} l'image de $y'$, et considérons
le diagramme de trois carrés
cartésiens ci-contre. Le lemme 7.5.
implique que $\,{}_{y'\backslash}Y'\longrightarrow{}_{y\backslash}Y'$ est l'inclusion
dans $\,{}_{y\backslash}Y'$ d'une composante connexe de $\,{}_{y\backslash}Y'$,
c'est donc une imm. ouverte. Donc
$\,{}_{y'\backslash}X\unsure{-|\to}{}_{y\backslash}X$ est une imm. ouverte. Comme
$\,{}_{y\backslash}X$ est loc. cofiltrante (resp loc fbt cof.) par
hypothèse, il en est donc de $=$ de
$\,{}_{y'\backslash}X$. Cela prouve a).

Appliquant a) à la situation $X_{/x}\longrightarrow Y_{/y}$
de b), on voit que les coloc\unsure{alisés}
$\,{}_{z\backslash}(X_{/x})$ de $X_{/x}$ sur $Y_{/y}$ $(z\in\operatorname{Ob}(Y_{/y}))$
sont loc. cofiltrants (resp fbt). Mais
comme $X_{/x}$ a un objet final\inserted{ au dessus de l'obj final de $Y_{/y}$}, il s'ensuit
que $\,{}_{z\backslash}(X_{/x})$ a un objet final, donc est \uline{\unsure{$0$-connexe}},
\margin{! $\vert\!\vert\!\vert$}
% archive sheet number: 239
\scanpage{240}\mspage{95}
Donc ces catégories, étant loc cofiltrantes (resp fbt),
sont en fait cofiltrantes (resp fbt cofiltrantes).

Prouvons c). Notons que pour $z$ dans $V$,
$\,{}_{z\backslash}U$ s'identifie à l'image inverse, par $\,{}_{z\backslash}(X_{/x})\unsure{-|\to}X_{/x}$,
de l'ouvert $U$ de $X_{/x}$, c'est donc un
\uline{ouvert} de $\,{}_{z\backslash}(X_{/x})$, non vide par construction
de $V$. Mais un ouvert non vide d'une
catégorie cofiltrante (resp. fbt. cofiltrante)
est encore cofiltrante (resp fbt. cofiltrante).
Cela signifie que $g:U\longrightarrow V$ est cofiltrant\deleted{\unsure{e}}
(resp fbt cofiltrant), a fortiori cofinal
(et $W$-coasphérique, pour tte $W$ satisfaisant
$\mathrm{Loc}(8)$). Cela achève la démonstration.

\uline{Remarque} \unsure{7.5}.\inserted{1)} Le même argument que dans
c) montre que si $f:X\unsure{-|\to}Y$ est
cofiltrant\unsure{e} (resp. fbt. cofiltrant), alors
pour \unsure{tout|tt} ouvert $U$ de $X$, si on
désigne par $V$ l'ouvert de $Y$ \deleted{\unsure{formé}} engendré
par $\unsure{f(U)}$, i.e. l'ouvert formé des $y$
de $Y$ tels que $\,{}_{y\backslash}X\ne\varnothing$, alors la
flèche induite $g:U\unsure{-|\to}V$ est cofiltrante
(resp. fbt cofiltrante), et \deleted{a fortiori}
est cofinal (et plus généralement, \inserted{si $f$ est cofiltrant, $g$ est}
$W$-coasphérique pour \deleted{\unsure{tout}}\inserted{tte} $W$ satisfaisant
\inserted{la condition} $\mathrm{Loc}(8)$).
% archive sheet number: 240
\scanpage{241}\mspage{96}
2) Un cas particulièrement intéressant \inserted{de 7.4}
est celui où on suppose que les $\,{}_{y\backslash}X$ $(y\in\operatorname{Ob}Y)$
sont \deleted{\unsure{toutes}} sommes directes de catégories
ayant chacune un \uline{objet initial}. Alors
on voit que cette condition est stable par
les opérations de 7.4. a), et \deleted{\unsure{que}}\unsure{par}
conséquent, les colocalisés \inserted{$\,{}_{z\backslash}(X_{/x})$} des morphismes
localisés en haut $X_{/x}\unsure{-|\to}Y_{/y}$ sont des
catégories à objet initial. Ceci \deleted{\unsure{reste}
\unsure{satisfait}} implique alors qu'il en est de
$=$ des colocalisés $\,{}_{z\backslash}U$ $(z\in\operatorname{Ob}V)$ dans
$g:U\unsure{-|\to}V$. Il s'ensuit que $\unsure{S|g}$ est
$W$-coasphérique pour \uline{tt} \deleted{\unsure{loc}}\unsure{localiseur} \unsure{''}
fondamental $W$ (sans avoir à se servir
de l'axiome des limites).

La condition envisagée sur $f:X\unsure{-|\to}Y$
est satisfaite notamment dans les
deux cas suivants (qui me paraissent
de loin les plus importants) : a) $f$ est
étale et b) $X$ est discrète.
% archive sheet number: 241
